% Kaelix — práticas de sistemas embarcados críticos e redundância aviônica
% Compilar:  latexmk -pdf praticas-sistemas-criticos.tex
%
% Documento de referência. Formato tabular deliberado: cada linha é um
% mecanismo, e as colunas fixam origem normativa, falha coberta, custo,
% limite conhecido e situação no Kaelix.

\documentclass[10pt,a4paper,landscape]{article}

\usepackage[T1]{fontenc}
\usepackage[utf8]{inputenc}
\usepackage[brazil]{babel}
\usepackage{lmodern}
\usepackage[landscape,top=1.6cm,bottom=1.7cm,left=1.5cm,right=1.5cm]{geometry}
\usepackage{booktabs}
\usepackage{tabularx}
\usepackage{array}
\usepackage{siunitx}
\usepackage{amsmath}
\usepackage{microtype}
\usepackage[hidelinks]{hyperref}
\usepackage{caption}
\usepackage{enumitem}
\usepackage{titlesec}
\usepackage{float}

\sisetup{output-decimal-marker={,}, group-separator={.}, per-mode=symbol}
\captionsetup{font=footnotesize, labelfont=bf}
\setlist{itemsep=1pt, topsep=3pt, parsep=0pt, leftmargin=*}
\titleformat{\section}{\normalfont\large\bfseries}{\thesection}{0.7em}{}
\titlespacing{\section}{0pt}{10pt}{5pt}
\pagestyle{plain}

% Colunas proporcionais: os pesos somam o número de colunas da tabela.
\newcolumntype{Y}[1]{>{\raggedright\arraybackslash\hsize=#1\hsize}X}

\begin{document}

\begin{center}
{\Large\bfseries Práticas de tolerância a falhas em sistemas embarcados críticos}\\[2pt]
{\large Redundância aviônica, normas de garantia e aplicabilidade ao Kaelix}\\[4pt]
{\footnotesize Documento de referência --- projeto Kaelix, Engenharia de Software (ICEV)}
\end{center}

\vspace{2pt}

\section{Escopo e critério de leitura}

\begin{itemize}
\item Unidade da tabela: \emph{mecanismo}, não norma. Uma norma aparece como origem de vários mecanismos.
\item A coluna \textbf{Limite conhecido} registra o que o mecanismo comprovadamente \emph{não} cobre. Sem ela, a tabela viraria lista de recomendações.
\item A coluna \textbf{Situação no Kaelix} admite quatro valores: \textsf{adotado}, \textsf{parcial}, \textsf{ausente}, \textsf{n/a} (não aplicável à classe do sistema).
\item Aviônica civil e automotivo entram como origem documentada de prática, não como alvo de certificação. O Kaelix é protótipo acadêmico; nenhum processo de certificação é reivindicado.
\end{itemize}

\section{Verificação das fontes}

\begin{itemize}
\item Conferidos em fonte primária: Knight e Leveson~\cite{knight1986} (cabeçalho, resumo, 27 versões, \num{1000000} testes); relatório da Comissão de Inquérito do Ariane 501~\cite{esa1996} (datas e causa).
\item Conferidos em registro de editor: Yeh~\cite{yeh1996} (veículo, volume, páginas); Brière e Traverse~\cite{briere1993} (veículo, páginas, DOI).
\item Normas (\cite{do178c,iec61508,iso26262,arinc653,misra2008,jsf2005}) citadas por edição e ano; texto integral não consultado. Numeração de regra conferida contra o uso no código do projeto.
\end{itemize}

\begin{table}[!p]
\centering
\caption{Mecanismos de tolerância a falhas em sistemas embarcados críticos. Origem, cobertura, custo, limite e situação no Kaelix.}
\label{tab:praticas}
\scriptsize
\renewcommand{\arraystretch}{1.25}
\begin{tabularx}{\textwidth}{@{}Y{0.80}Y{0.80}Y{1.10}Y{0.82}Y{1.26}Y{1.22}@{}}
\toprule
\textbf{Mecanismo} & \textbf{Origem} & \textbf{Falha coberta} & \textbf{Custo} & \textbf{Limite conhecido} & \textbf{Situação no Kaelix} \\
\midrule

Redundância modular tripla com votação &
Boeing 777, três canais de mesmo número de peça~\cite{yeh1996} &
Falha aleatória de hardware; transiente em um canal &
3$\times$ canais; sincronização de quadro &
Nula contra defeito de projeto replicado nos três canais &
\textsf{ausente} --- nó único \\

Dissimilaridade de projeto ($N$-versões) &
777: 3 microprocessadores e 3 compiladores Ada~\cite{yeh1996}; A320: ELAC/SEC de fabricantes distintos~\cite{briere1993}; Shuttle: PASS e BFS por empresas distintas~\cite{nasa2011} &
Modo comum de projeto, de compilador e de cadeia de ferramentas &
2--3$\times$ esforço de desenvolvimento &
Independência não se confirma: 27 versões da mesma especificação, \num{1000000} testes, falhas correlacionadas acima do esperado~\cite{knight1986} &
\textsf{parcial} --- paridade numérica C++\,$\leftrightarrow$\,Python (\emph{Isolation Forest}, CRC-16) em teste; sem votação em operação \\

Par comando/monitor &
A320, canal duplo dissimilar por computador~\cite{briere1993} &
Saída errada aceita como válida (\emph{runaway} de superfície) &
2$\times$ canal por unidade &
Monitor herda a especificação do comando: erro de especificação passa nos dois &
\textsf{ausente} --- veredito de anomalia sem segunda opinião \\

Votação por valor mediano &
777, seleção entre três pistas~\cite{yeh1996} &
Divergência numérica entre réplicas &
Seleção por quadro de dados &
Exige três fontes ou mais; mediana de duas é indefinida &
\textsf{ausente} \\

Watchdog independente, em janela ou externo &
IEC 61508~\cite{iec61508}; ISO 26262~\cite{iso26262} &
Travamento; laço de realimentação curto demais; base de tempo parada &
Temporizador dedicado; pino e circuito integrado externo &
Não detecta resultado errado produzido no prazo certo; watchdog que compartilha relógio com o núcleo compartilha a falha &
\textsf{parcial} --- WDT de tarefa e cota de software ativos; WDT de domínio RTC não armável no ESP32-S3 (\texttt{soc/rtc\_wdt.h} referencia registradores inexistentes no alvo); sem watchdog externo \\

Partição temporal e espacial &
ARINC 653, APEX~\cite{arinc653} &
Propagação de falha entre funções de criticidade distinta no mesmo processador &
RTOS particionado; orçamento fixo de tempo e memória por partição &
Não se aplica a executivo de função única &
\textsf{n/a} --- \emph{superloop} de um disparo, uma função \\

\emph{Lockstep} de núcleo duplo com comparador &
ISO 26262, mecanismo dominante para ASIL D~\cite{iso26262} &
Falha transiente de núcleo, detectada ciclo a ciclo &
2$\times$ área de núcleo; sem ganho de desempenho &
Réplicas idênticas executam o mesmo defeito de software simultaneamente &
\textsf{ausente} --- ESP32-S3 não oferece \emph{lockstep} \\

Estado seguro definido e alcançável &
IEC 61508~\cite{iec61508} &
Comportamento indefinido após detecção de falha &
Projeto do caminho de saída; verificação de que é alcançável de qualquer ponto &
Pressupõe que exista estado seguro; sistemas \emph{fail-operational} não podem apenas desligar &
\textsf{adotado} --- sequência de três passos (rádio dormindo, periféricos cortados, sono armado) com escalada de período após resets anormais consecutivos \\

Memória estática, sem alocação dinâmica &
JSF++ AV-206~\cite{jsf2005}; MISRA C++ 18-4-1~\cite{misra2008} &
Exaustão e fragmentação de \emph{heap}; falha de alocação não tratada &
Pico de RAM dimensionado em tempo de compilação &
Desloca o limite para o projeto: dimensionamento errado falha no build, não em campo &
\textsf{adotado} --- buffers de amostras e de FFT em \texttt{.bss} \\

Norma de codificação com verificação automática &
MISRA C++:2008~\cite{misra2008}; JSF++~\cite{jsf2005} &
Construção ambígua; comportamento indefinido; conversão implícita &
Analisador estático integrado ao build &
Conformidade não implica ausência de defeito nem correção funcional &
\textsf{adotado} --- norma própria derivada, \texttt{clang-tidy}, \texttt{clang-format}, \texttt{-Werror} \\

Verificação de invariante em tempo de compilação &
Prática geral de sistemas críticos; coerente com MISRA C++~\cite{misra2008} &
Divergência silenciosa de layout, tamanho ou contrato entre módulos &
Nulo em tempo de execução &
Só alcança o que é constante em compilação; exige que o build do alvo seja executado &
\textsf{adotado} --- asserções de tamanho de quadro, deslocamento de campo e prazo do ciclo, avaliadas pela primeira vez em build de alvo \\

Rastreabilidade requisito $\rightarrow$ código $\rightarrow$ teste &
DO-178C, tabelas do Anexo A~\cite{do178c} &
Requisito implementado sem verificação correspondente &
Documental e contínuo &
Rastreia cobertura, não correção; matriz desatualizada descreve sistema que não existe &
\textsf{adotado} --- matriz declara explicitamente o não verificado e o motivo \\

Verificação com independência &
DO-178C: objetivos a cumprir com independência nos níveis mais altos~\cite{do178c} &
Viés do autor ao verificar o próprio artefato &
Pessoal distinto de quem produziu o item; DAL A estimado em cerca de 3$\times$ o custo de DAL B/C &
Independência organizacional não elimina erro de especificação compartilhada &
\textsf{ausente} --- autor único \\

Detecção de erro fim a fim por código detector &
Prática geral em enlaces sem confirmação &
Corrupção silenciosa em canal de rádio e em memória retida &
2 bytes por quadro; ciclos de cálculo &
Detecta, não corrige; não cobre erro cometido antes do cálculo do código &
\textsf{adotado} --- CRC-16/CCITT no quadro de rádio e no bloco retido em memória RTC \\

Reanálise de envelope em reuso de software &
Ariane 501: unidade inercial herdada do Ariane 4~\cite{esa1996} &
Premissa de faixa numérica válida no sistema de origem e inválida no novo &
Reanálise por item reusado &
As duas unidades redundantes falharam pela mesma causa em segundos: redundância por réplica não cobre premissa herdada &
\textsf{risco aberto} --- limiares de norma de vibração e correntes de folha de dados adotados sem reanálise no envelope do projeto \\

\bottomrule
\end{tabularx}
\end{table}


\section{Escalas de criticidade}

\begin{table}[H]
\centering
\caption{Escalas de criticidade por setor. Não há equivalência normativa entre elas; a correspondência citada na literatura é informal.}
\label{tab:escalas}
\footnotesize
\renewcommand{\arraystretch}{1.2}
\begin{tabularx}{\textwidth}{@{}Y{0.72}Y{0.55}Y{1.25}Y{1.48}@{}}
\toprule
\textbf{Norma (setor)} & \textbf{Níveis} & \textbf{Base de atribuição} & \textbf{Observação} \\
\midrule
DO-178C (aviônica, software)~\cite{do178c} & A--E & Severidade da condição de falha, de catastrófica a sem efeito de segurança & Não atribui taxa de falha a software; as probabilidades vêm da avaliação de segurança do sistema, não deste documento \\
IEC 61508 (genérica)~\cite{iec61508} & SIL 1--4 & Medida de falha alvo, por hora ou por demanda & Única das três com faixa numérica direta (Tab.~\ref{tab:sil}) \\
ISO 26262 (automotivo)~\cite{iso26262} & QM, ASIL A--D & Severidade $\times$ exposição $\times$ controlabilidade & QM indica ausência de requisito de segurança funcional \\
\bottomrule
\end{tabularx}
\end{table}

\begin{table}[H]
\centering
\caption{IEC 61508: medida de falha alvo por nível de integridade~\cite{iec61508}.}
\label{tab:sil}
\footnotesize
\renewcommand{\arraystretch}{1.2}
\begin{tabular}{@{}lll@{}}
\toprule
\textbf{Nível} & \textbf{Alta demanda ou contínuo (por hora)} & \textbf{Baixa demanda (por demanda)} \\
\midrule
SIL 4 & $\geq \num{1e-9}$ a $< \num{1e-8}$ & $\geq \num{1e-5}$ a $< \num{1e-4}$ \\
SIL 3 & $\geq \num{1e-8}$ a $< \num{1e-7}$ & $\geq \num{1e-4}$ a $< \num{1e-3}$ \\
SIL 2 & $\geq \num{1e-7}$ a $< \num{1e-6}$ & $\geq \num{1e-3}$ a $< \num{1e-2}$ \\
SIL 1 & $\geq \num{1e-6}$ a $< \num{1e-5}$ & $\geq \num{1e-2}$ a $< \num{1e-1}$ \\
\bottomrule
\end{tabular}
\end{table}

\section{Leitura da coluna de situação}

\begin{itemize}
\item Sete mecanismos \textsf{adotados} concentram-se em disciplina de codificação, detecção de erro e definição de estado seguro --- categorias cujo custo é de projeto, não de replicação de hardware.
\item Os mecanismos \textsf{ausentes} são, sem exceção, os que exigem canal, núcleo ou unidade redundante. A ausência é consequência direta da topologia de nó único alimentado a bateria, não de omissão de projeto.
\item Dois itens exigem ação: o watchdog de domínio RTC, indisponível no alvo pela via de API adotada; e o reuso de limiares e de valores de folha de dados sem reanálise de envelope, cuja consequência histórica está documentada em~\cite{esa1996}.
\item A dissimilaridade parcial do Kaelix é de \emph{verificação}, não de \emph{operação}: duas implementações independentes são comparadas em teste, e apenas uma executa em campo. O ganho é contra erro de transcrição entre treino e inferência, não contra defeito de especificação --- distinção que~\cite{knight1986} torna obrigatória.
\end{itemize}

\begin{thebibliography}{99}
\footnotesize

\bibitem{arinc653}
AERONAUTICAL RADIO, INC. \emph{ARINC Specification 653: Avionics Application Software Standard Interface}. Annapolis: ARINC, 1996. Suplemento 1 (1997) introduz o APEX e o particionamento temporal e espacial.

\bibitem{briere1993}
BRIÈRE, D.; TRAVERSE, P. AIRBUS A320/A330/A340 electrical flight controls --- a family of fault-tolerant systems. In: \emph{FTCS-23: The Twenty-Third International Symposium on Fault-Tolerant Computing}. Toulouse, 22--24 jun. 1993. p.~616--623. DOI: 10.1109/FTCS.1993.627364.

\bibitem{esa1996}
AGÊNCIA ESPACIAL EUROPEIA. \emph{Ariane 501 --- Presentation of Inquiry Board report}. Paris, 23 jul. 1996. Falha em 4 jun. 1996: perda completa de informação de guiagem e atitude 37\,s após o início da sequência de ignição, com as duas unidades inerciais perdidas pela mesma causa. Disponível em: \url{https://www.esa.int/Newsroom/Press_Releases/Ariane_501_-_Presentation_of_Inquiry_Board_report}.

\bibitem{iec61508}
INTERNATIONAL ELECTROTECHNICAL COMMISSION. \emph{IEC 61508: Functional safety of electrical/electronic/programmable electronic safety-related systems}. 2.~ed. Genebra: IEC, 2010.

\bibitem{iso26262}
INTERNATIONAL ORGANIZATION FOR STANDARDIZATION. \emph{ISO 26262: Road vehicles --- Functional safety}. 2.~ed. Genebra: ISO, 2018.

\bibitem{jsf2005}
LOCKHEED MARTIN CORPORATION. \emph{Joint Strike Fighter Air Vehicle C++ Coding Standards for the System Development and Demonstration Program}. Doc. 2RDU00001 Rev.~C, dez. 2005.

\bibitem{knight1986}
KNIGHT, J. C.; LEVESON, N. G. An experimental evaluation of the assumption of independence in multiversion programming. \emph{IEEE Transactions on Software Engineering}, v.~SE-12, n.~1, p.~96--109, jan. 1986. DOI: 10.1109/TSE.1986.6312924.

\bibitem{misra2008}
MIRA LIMITED. \emph{MISRA C++:2008 --- Guidelines for the use of the C++ language in critical systems}. Nuneaton: MIRA, 2008.

\bibitem{nasa2011}
NATIONAL AERONAUTICS AND SPACE ADMINISTRATION. \emph{Space Shuttle GN\&C Development History and Evolution}. NTRS 20110014833. Houston: NASA, 2011. Sistema primário com quatro computadores em conjunto redundante e um quinto executando software de reserva desenvolvido separadamente.

\bibitem{do178c}
RTCA, INC. \emph{DO-178C: Software Considerations in Airborne Systems and Equipment Certification}. Washington: RTCA, dez. 2011.

\bibitem{yeh1996}
YEH, Y. C. Triple-triple redundant 777 primary flight computer. In: \emph{Proceedings of the 1996 IEEE Aerospace Applications Conference}. IEEE, 1996. v.~1, p.~293--307.

\end{thebibliography}

\end{document}
